\documentclass[10pt,a4paper]{amsart}
\usepackage[margin=19mm]{geometry}
\usepackage{amsmath,amssymb,mathtools}
\usepackage[scr=rsfs,cal=euler]{mathalfa}
\usepackage{xfrac}
\usepackage[unicode,hidelinks,bookmarks=false]{hyperref}
\newcommand{\ie}{i.e.\ }
\newcommand{\cf}{cf.\ }
\newcommand{\resp}{respectively\ }
\newcommand{\Subsection}{\subsection}
\providecommand{\nobreakdash}{\nobreak\hbox{-}}
\setlength{\emergencystretch}{2em}
\multlinegap0pt
\usepackage{amssymb}
\usepackage{stackengine}
\newtheorem{prop}{Proposition}
\newcommand{\D}{\mathcal{D}}
\def\bbbr{\mathbb{R}}
\newcommand\lbd{\underline{\mbox{\rm dim}}_{\rm B}} % lower box dimension
\newcommand\mlbd{\underline{\mbox{\rm dim}}_{\rm MB}} %modified lower box
\newcommand\pkd{\mbox{\rm dim}_{\rm P}} % packing dimension
\newcommand\ubd{\overline{\mbox{\rm dim}}_{\rm B}} % upper box dimension
\newcommand\ucod{\overline{\mbox{\rm dim}}_{\rm C}} % upper correlation dimension
\title{Frostman and Fourier characterisations of fractal dimensions}
\author{Kenneth J. Falconer, Shuqin Zhang}
\date{Source: arXiv:2505.21217v2 (CC BY 4.0)}
\begin{document}
\maketitle

\noindent\emph{Source and licence.} Frostman and Fourier characterisations of fractal dimensions. Version arXiv:2505.21217v2 (CC BY 4.0), distributed by arXiv under the Creative Commons Attribution 4.0 International licence, http://creativecommons.org/licenses/by/4.0/, as recorded in the arXiv licence metadata for this exact version and shown on the abstract page https://arxiv.org/abs/2505.21217v2. Full licence notice: \texttt{licenses/arxiv-2505.21217-CC-BY-4.0.txt}.\par\smallskip

\noindent\textit{Editorial scope.} Counted unit: the article's prop example1 (source0.tex lines 661--663). The article's complete original proof of it is delivered with it (source0.tex lines 665--809, one \texttt{proof} environment of 145 lines). No mathematics is written, repaired or re-derived: the statement and the proof are the article's own lines, in the article's own notation, and no retained line is substituted or re-flowed.  Retained uncounted as the source-local context the proof needs: lines 286--295.  Nothing was omitted from any retained span, so the package carries no omission record.  No retained line contains a citation, so the wrapper carries no bibliography and no bibliography entry.  No retained line contains a figure environment, an image-inclusion command or drawing code, so no image is delivered and the package declares no asset dependency.  Editorial formatting only, in addition: an AMS \texttt{amsart} preamble that restates the article's own declarations and macros used by the retained lines, namely the environments \texttt{prop} on separate counters, together with the standard packages those lines need.  The retained environments print in this document as Proposition 1 in order of appearance, whereas the article prints the same items with the numbering of its own full document.  The complete native source of the article as distributed in the arXiv e-print for this exact version is bundled unchanged as \texttt{sources/178850/source0.tex}.
\begin{prop}\label{cordim}
Let $E \subset \mathbb{R}^d$ be a Borel set. Then
\begin{align}
&\ucod E\nonumber\\
&= \inf\Big\{s \geq 0:  \forall\ \mu \in {\mathcal P}(E), \text{ for all sufficiently small }r, \int \mu(B(x,r))\mathrm{d}\mu(x) \geq r^{s } \Big\} \label{cd1}\\
&= \sup\Big\{s \geq 0: \exists\ \mu \in {\mathcal P}(E)  \text{ and $\{r_k\}_k\searrow 0 $ s.t.  } \int \mu(B(x,r_k))\mathrm{d}\mu(x) \leq r_k^{s } \Big\} \label{cd2} \\
&= \sup\Big\{s \geq 0: \exists\   \mu \in {\mathcal P}(E)  \text{ and $\{r_k\}_k\searrow 0 $ s.t.  } \mu(B(x,r_k))\leq r_k^{s } \text{ for all } x \in \mathbb{R}^d\Big\}\label{cd3}\\
&= \sup\Big\{s \geq 0: \exists\   \mu \in {\mathcal P}(E)  \text{ and $\{n_k\}_k\nearrow\infty $ s.t.  } \mu(C_{n_k})\leq 2^{-n_ks}  \text{ for all } C_{n_k}  \in \D_{n_k}\Big\}\label{cd4}.
\end{align}
\end{prop}

\begin{prop}\label{example1}
Given $0<t<s<1$, there exists a compact set $E\subset \bbbr^n$ such that $\mlbd E=\lbd E  = t$  and $\ucod E=\pkd E =s$.
\end{prop}

	\begin{proof}
		Fix positive real numbers $0<t<s<1$. Let $\{\epsilon_k\}_k\subset (0,t)$ be a sequence of real numbers  decreasing to $0$. Let  $n_0=0$ and $n_1=1$.  We define a sequence of integers $\{n_m\}_m$ recursively.
		 For $k\geq1$, let
		\begin{equation}\label{n2k}
			n_{2k}=\bigg\lfloor\frac{1}{t-\epsilon_k}\sum_{i=0}^{k-1}(n_{2i+1}-n_{2i})  \bigg\rfloor+1
		\end{equation}	
		and 
		\begin{equation}\label{n2k1}
			n_{2k+1}=\bigg\lfloor\frac{1}{1-s}\Big(n_{2k}-\sum_{i=0}^{k-1}(n_{2i+1}-n_{2i})\Big)\bigg\rfloor,
		\end{equation}	
	 where $\lfloor c\rfloor$ is  the integral part of the real number $c$.
		 Define $\mathcal{D}_n$ as the collection of all half-open dyadic intervals of length  $2^{-n}$  that are contained in $[0,1]$, which we refer to as $n$th-level intervals, so 
		\begin{equation*}
			\mathcal{D}_n=\bigg\{\Big(\frac{k}{2^n},\frac{k+1}{2^n}\Big]:k=0,1,\cdots,2^n-1\bigg\}.
		\end{equation*}

Beginning with $\mathcal{C}_0=[0,1]$, for all   $n\in\mathbb{N}$ we will select dyadic intervals from $\mathcal{D}_n$, and  denote the set of all such chosen intervals as $\mathcal{C}_n$. 
		For each $n$th-level interval chosen for $\mathcal{C}_n$,  we ensure that at least one of  its $(n+1)$th-level subintervals will be chosen for $\mathcal{C}_{n+1}$. 
		 Thus at the $(n+1)$th stage, we either choose the lefthand $(n+1)$th-level subinterval of  each interval in $\mathcal{C}_n$ (referred to as {\it Method 1})  or we choose both	$(n+1)$th-level subintervals of  each  interval in $\mathcal{C}_n$ (referred to as  {\it Method 2}).
		 
		At the first stage ($n_1=1$), we adopt Method 2. Thus
		\begin{equation*}
			\textstyle{\mathcal{C}_{n_1}=\big\{(0,\frac{1}{2}],(\frac{1}{2},1]\big\}.}
		\end{equation*}
		We alternate between using Method 1 and Method 2 throughout the construction process.
		At the $m$th stage, we use Method 1 when $n_{2k-1}<m\leq n_{2k}$ for an integer $k$,  and Method 2 when $n_{2k}<m\leq n_{2k+1}$.
			Define
		\begin{equation*}
			E_n=\bigcup_{C_n\in \mathcal{C}_n}C_n \quad\textit{and}\quad E=\bigcap_{n =1}^\infty \overline{E_n}\, ,
		\end{equation*}
		where $\overline{E_n}$ is the closure of $E_n$.
		
		Clearly, when $n_{2k-1}\leq m\leq n_{2k}$,
		\begin{equation*}
			\#(\mathcal{C}_{n_{2k}}\cap C_{m})=\left\{
			\begin{aligned}
				&1 \quad&\text{if } C_{m}\in\mathcal{C}_{m};\\
				&0 \quad&\text{if } C_{m}\notin\mathcal{C}_{m}.
			\end{aligned}\right.
		\end{equation*}
		Thus 
		\begin{equation*}
			\#\mathcal{C}_{n_{2k}}=\prod_{i=0}^{k-1}2^{n_{2i+1}-n_{2i}}.
		\end{equation*}
		By \eqref{n2k}, 
		\begin{equation}\label{2k}
			2^{n_{2k}(t-\epsilon_k)-1}\leq\#\mathcal{C}_{n_{2k}}<2^{n_{2k}(t-\epsilon_k)}
		\end{equation}
		When $n_{2k}\leq m\leq n_{2k+1}$,
		\begin{equation*}
			\#(\mathcal{C}_{n_{2k+1}}\cap C_{m})=\left\{ 
			\begin{aligned}
				&2^{n_{2k+1}-m} \quad&\text{if } C_{m}\in\mathcal{C}_{m};\\
				&0 \quad&\text{if } C_{m}\notin\mathcal{C}_{m}.
			\end{aligned}\right.
		\end{equation*}
		Consequently,
		\begin{equation*}
			\#\mathcal{C}_{n_{2k+1}}=\prod_{i=0}^{k}2^{n_{2i+1}-n_{2i}}.
		\end{equation*}
		It follows from \eqref{n2k1} that
		\begin{equation}\label{2k+1}
			2^{n_{2k+1}s-1}<\#\mathcal{C}_{n_{2k+1}}\leq2^{n_{2k+1}s}.
		\end{equation}
		
		It follows directly from \eqref{2k} that $\lbd E\leq t$.
		We next show that $\mlbd E\geq t$.
		Let $\epsilon>0$ and
		let $E=\bigcup_i E_i$ be an arbitrary countable cover of $E$ with each  $E_i$  closed.
		Since $E$ is compact, by Baire's category theorem there  exists an  $E_i$ such that $E \cap E_i$ has non-empty interior relative to $E$.
		Thus, for some $n$, there exists a $n$th-level interval $C_n$ such that $E\cap C_n\subset  E_i$.
		Since $E\cap C_n\neq\emptyset$, it follows that for $m>n$,
		\begin{align*}
			N_{2^{-m}}(E_i)&\geq N_{2^{-m}}(E\cap C_n)\\
			&\geq\frac{\#\mathcal{C}_m}{\#\mathcal{C}_n}\\
			&=\frac{1}{\#\mathcal{C}_n}\begin{cases}
				\#\mathcal{C}_{n_{2k}}, & \text{if } n_{2k-1}<m\leq n_{2k}\\
					\#\mathcal{C}_{n_{2k}}\cdot 2^{m-n_{2k}}, & \text{if } n_{2k}<m\leq n_{2k+1}
			\end{cases}\\
			&\geq\frac{1}{\#\mathcal{C}_n}\begin{cases}
				2^{n_{2k}(t-\epsilon_k)-1}, & \text{if } n_{2k-1}<m\leq n_{2k}\\
				2^{n_{2k}(t-\epsilon_k)-1+(m-n_{2k})}, & \text{if } n_{2k}<m\leq n_{2k+1}
			\end{cases}\quad\big(\text{by } \eqref{2k}\big)\\
			&\geq\frac{1}{2\#\mathcal{C}_n}2^{m(t-\epsilon)},
		\end{align*}
		where the last inequality is valid if $m$ is large enough so that the corresponding $k$ satisfies $\epsilon_k<\epsilon$.
		Thus, $\lbd E_i\geq t-\epsilon$. According to the definition of the modified lower box dimension and the arbitrariness of $\epsilon$, we conclude $\mlbd E\geq t$
		and thus
		\begin{equation*}
			\mlbd E=\lbd E= t.
		\end{equation*}
		
		
		
		We next  verify that  $\ucod (E)= \pkd E = s$. For each $m\in \mathbb{N}$, define a measure 
		\begin{equation*}
			\mu_m= \sum_{C\in\mathcal{C}_{n_{2m+1}}}\frac{1}{\#\mathcal{C}_{n_{2m+1}}}\frac{\mathcal{L}|_C}{\mathcal{L}(C)},
		\end{equation*}
		where $\mathcal{L}$ is Lebesgue measure.
		For  $C_{n_{2k+1}}\in\mathcal{C}_{n_{2k+1}}$ and $k\leq m$, 
			\begin{align}\label{mumc}
			\mu_m(C_{n_{2k+1}})
			=&\ \#(\mathcal{C}_{n_{2m+1}}\cap C_{n_{2k+1}})\cdot\frac{1}{\#\mathcal{C}_{n_{2m+1}}}\notag\\
			=&\ \frac{1}{\#\mathcal{C}_{n_{2k+1}}}\notag\\
			\leq&\ 2^{-n_{2k+1}s+1}\quad\big(\text{by } \eqref{2k+1}\big).
		\end{align}	
		
		Let $\{\mu_{m_j}\}_j$ be a weakly convergent subsequence of $\{\mu_m\}_m$, and let $\mu$ be the limit of this subsequence.
		It is easy to see that $\mu\in\mathcal{P}(E)$.
		For each $x\in E$, $x$ lies in the interior of the union of two ${n_{2k+1}}$th adjacent dyadic intervals, denoted as $U_k(x)$.
			Thus
		\begin{align*}
			\mu(\{x\})\leq&\ \mu(U_k(x))\\
		\leq&\ \liminf_m\mu_m(U_k(x))\\
		\leq&\  2^{-n_{2k+1}s+2}\quad\big(\text{by } \eqref{mumc}\big).
		\end{align*}	
			Letting  $k$ tend to infinity  implies that $\mu$ has no atoms.
			We denote the interior of  ${C}_{n_{2k+1}}$ by ${\rm int }\,{C}_{n_{2k+1}}$ and its boundary by  $\partial C_{n_{2k+1}}$. Since $\mu(\partial C_{n_{2k+1}})=0$, 
		\begin{align*}
			\mu(C_{n_{2k+1}})\leq&\ \mu(\partial C_{n_{2k+1}})+\mu({\rm int }\,{C}_{n_{2k+1}})\\
			\leq&\ \liminf_m\mu_{m}({\rm int }\,{C}_{n_{2k+1}})\\
			\leq&\ 2^{-n_{2k+1}s+1}\quad\big(\text{by } \eqref{mumc}\big).
		\end{align*}	
By Proposition \ref{cordim},  $\ucod E\geq s$.
		
		If $\pkd  E> s$, then there exists $\epsilon_0>0$ such that $\ubd(E)\geq\pkd E> s+\epsilon_0$.
		Then  there exists a sequence of integers $\{m_k\}_k$ such that
		\begin{equation}\label{contra}
			\#\mathcal{C}_{m_k}\geq 2^{m_k(s+\epsilon_0)}.
		\end{equation}
		However, for each $m\in\mathbb{N}$, 
			\begin{align*}
			\#\mathcal{C}_m&=\begin{cases}
				\#\mathcal{C}_{n_{2k-1}}, & \text{if } n_{2k-1}<m\leq n_{2k}\\
				\#\mathcal{C}_{n_{2k+1}}\cdot 2^{m-n_{2k+1}}, & \text{if } n_{2k}<m\leq n_{2k+1}
			\end{cases}\\
			&\leq\begin{cases}
				2^{n_{2k-1}s}, & \text{if } n_{2k-1}<m\leq n_{2k}\\
				2^{n_{2k+1}s+m-n_{2k+1}}, & \text{if } n_{2k}<m\leq n_{2k+1}
			\end{cases}\quad\big(\text{by } \eqref{2k+1}\big)\\
			&\leq 2^{ms},
		\end{align*}
		which contradicts  \eqref{contra}.
		Thus $$\ucod E= \pkd E =s.$$
	\end{proof}

\end{document}
